\textbf{Spatiotemporal Action Localization.}
This task (STAL) is to predict the frames and corresponding actions of people in video keyframes.
We evaluate \model~on two classic STAL datasets AVA2.2~\cite{ava} and AVA-Kinetics~\cite{avakinetics}. In AVA2.2~\cite{ava}, each video lasts 15 minutes, and it gives a keyframe every second. The annotation is provided for keyframes instead of all frames. Here we use a classic two-stage approach to handle this task. We apply a well-trained (on MS-COCO \cite{lin2014microsoft}) Mask-RCNN \cite{maskrcnn} to detect humans on each keyframe, and the keyframe boxes are provided in the Alphaction~\cite{alphaction} project. In the second stage, centering around the key frame, a certain number of frames are extracted and fed into our video backbone.
Similarly, in the training, we use the ground truth box for training~\cite{alphaction}, and the boxes predicted in the first stage for testing. 


We used ViT-Huge in \model~for experiments. The specific results can be seen in Table \ref{tab:stal}. The classification head uses a simple linear head, achieving the SOTA performance on both datasets. Note using the ViT-H model, and training with the AVA-Kinetics dataset not only improves the overall mAP, but also significantly improves the mAP obtained by testing on AVA alone. It suggests the introduction of some Kinetics videos to AVA will improve the generalization of the model over AVA; on the other hand, observing the various distributions of the AVA dataset, we find that AVA presents a typical long-tailed distribution. The introduction of Kinetics video will alleviate this issue for better results. Due to the small number of models validated on the AVA-Kinetics dataset, only the results from the \href{https://paperswithcode.com/sota/spatio-temporal-action-localization-on-ava}{paperswithcode} website are selected in the Table \ref{tab:stal}.